%%%%%%%% FMSD @ ICML 2026 SUBMISSION %%%%%%%%%%%%%%%%%%%%%%%%%%%%
% Multimodal Time-Series Probing -- Disentangling Text from Priors
%%%%%%%%%%%%%%%%%%%%%%%%%%%%%%%%%%%%%%%%%%%%%%%%%%%%%%%%%%%%%%%%

\documentclass{article}

\usepackage{microtype}
\usepackage{graphicx}
\usepackage{subcaption}
\usepackage{booktabs}
\usepackage{adjustbox}
\usepackage{multirow}
\usepackage{array}
\usepackage{xcolor}
\usepackage{hyperref}

% Blind-review version of the ICML 2026 style.
\usepackage{icml2026}

\usepackage{amsmath}
\usepackage{amssymb}
\usepackage{mathtools}
\usepackage{amsthm}
\usepackage[capitalize,noabbrev]{cleveref}

% --- Best / second-best highlight commands -----------------------------
% Used by the auto-generated tab_main.tex.
\newcommand{\best}[1]{\textbf{#1}}
\newcommand{\second}[1]{\underline{#1}}

% --- Condition shortcuts ----------------------------------------------
\newcommand{\Cone}{\textsc{C1}}
\newcommand{\Ctwo}{\textsc{C2}}
\newcommand{\Cthree}{\textsc{C3}}
\newcommand{\Cfour}{\textsc{C4}}
\newcommand{\Cfive}{\textsc{C5}}
\newcommand{\Csix}{\textsc{C6}}
\newcommand{\Ceight}{\textsc{C8}}
\newcommand{\Cnine}{\textsc{C9}}

\newcommand{\TaTS}{\textsc{TaTS}}
\newcommand{\MMTSF}{\textsc{MM-TSFlib}}
\newcommand{\Aurora}{\textsc{Aurora}}
\newcommand{\TimeMMD}{Time-MMD}

% --- Math notation shortcuts ------------------------------------------
% A small, consistent vocabulary used in every equation.
\newcommand{\xnum}{\mathbf{x}}              % numeric history
\newcommand{\Ttext}{\mathbf{T}}             % text history
\newcommand{\pprior}{\mathbf{p}}            % prior column
\newcommand{\yhat}{\hat{\mathbf{y}}}        % forecast
\newcommand{\ytarget}{\mathbf{y}}           % ground truth
\newcommand{\Etext}{\mathbf{E}}             % text covariate matrix
\newcommand{\Ztext}{\mathbf{Z}_{\!T}}       % distilled text representation
\newcommand{\Real}{\mathbb{R}}
\newcommand{\Text}{\mathcal{T}}
\newcommand{\bbE}{\mathbb{E}}

\icmltitlerunning{Reading or Reciting? Multimodal Time Series on \TimeMMD}

\begin{document}

\twocolumn[
  \icmltitle{What Are Multimodal Time-Series Models Actually Reading?\\
             Disentangling Text Semantics from Numeric Priors on Time-MMD}

  \icmlsetsymbol{equal}{*}
  \begin{icmlauthorlist}
    \icmlauthor{Anonymous Authors}{anon}
  \end{icmlauthorlist}
  \icmlaffiliation{anon}{Anonymous Institution}

  \icmlcorrespondingauthor{Anonymous Authors}{anon@anon.invalid}
  \icmlkeywords{multimodal time series, foundation models, probing,
                evaluation, Time-MMD, text-as-context}
  \vskip 0.3in
]

\printAffiliationsAndNotice{}

% =====================================================================
\begin{abstract}
Multimodal time-series forecasters now claim that natural-language
context substantially improves numeric prediction.  Three concurrent
paradigms report double-digit MSE gains over unimodal baselines on
the \TimeMMD{} benchmark: a pretrained foundation model (\Aurora{}),
an early-fusion text-channel approach (\TaTS{}), and a late-fusion
residual approach (\MMTSF{}).  We test whether these gains stem from
\emph{text semantics} or from a co-shipped numeric \emph{prior column}
that flows through the same fusion path.  Holding model checkpoints
fixed, we perturb \TimeMMD{} along two orthogonal axes -- text content
(empty, shuffled, cross-domain, constant, oracle) and numeric prior
(intact, zeroed) -- yielding eight conditions.  Across nine domains,
four horizons, three seeds, and three backbones per fusion
architecture (\textbf{$7{,}776$ forecasting runs}), we find:
(i) zeroing the prior column degrades MSE by up to $31.1\%$, while
every text-only perturbation moves MSE by less than $0.06\%$;
(ii) \Aurora{}'s predictions are bit-identical with full text, zeroed
text, and the model's own unimodal flag, even though its text encoder
still receives non-zero gradients; (iii) \Aurora{}'s text
cross-attention has near-uniform entropy ($\widehat H = 0.975$)
regardless of input.  On \TimeMMD{}, the apparent multimodal gain is
structural rather than semantic, with implications for how benchmarks
report multimodal lift and how foundation-model evaluation should be
designed.
\end{abstract}

% =====================================================================
\section{Introduction}
\label{sec:intro}

Foundation models that consume both numeric time series and natural
language are now a fast-growing branch of structured-data
research~\citep{liu2024timemmd,zhou2025tats,wu2025aurora,liu2024autotimes}.
On \TimeMMD{}~\citep{liu2024timemmd} -- the standard text--time-series
forecasting benchmark -- three architectural paradigms have emerged.
A zero-shot pretrained model, \Aurora{}~\citep{wu2025aurora}, attends
over a tokenised text column.  A per-domain trained early-fusion
model, \TaTS{}~\citep{zhou2025tats}, projects per-row text embeddings
and concatenates them as additional input channels.  A per-domain
trained late-fusion model, \MMTSF{}~\citep{liu2024timemmd}, blends a
pooled-text vector with the time-series head's output.  All three
publish double-digit MSE improvements over unimodal baselines.

A closer look at the data pipeline complicates this story.  Every
\TimeMMD{} domain ships not only paired text but also a numeric
column \texttt{prior\_history\_avg}: an LLM-generated forecast for
the same horizon~\citep{liu2024timemmd}.  Two of the three reference
implementations actively use it.  In \MMTSF{}'s released training
loop the multimodal output is computed as
\(\texttt{prompt\_y} = \texttt{norm}(\texttt{prompt\_emb})
                       + \texttt{prior\_y}\),
where \texttt{prompt\_emb} is a frozen-BERT pooled text embedding and
\texttt{prior\_y} is the prior column.  \TaTS{} adds an analogous
output-side blend
$\yhat = (1\!-\!w)\,\yhat_{\text{model}}
        + w\cdot\pprior_{L+1:L+H}$
on top of an early-fusion text channel.  \Aurora{} does not consume
the prior column, but does cross-attend over the text column.
Crucially, in every published implementation, ``no text'' is
operationalised in a way that simultaneously zeroes a numeric signal
that has nothing to do with text semantics -- a confound that has
previously caused well-known reasoning benchmarks to collapse under
adversarial probes~\citep{niven2019probing,mccoy2019hans}.

We therefore ask: \emph{when these models appear to ``use'' text, is
the gain coming from natural-language semantics, or from the prior
column that flows through the same fusion path?}  We construct an
$8$-condition behavioural ladder over two orthogonal axes -- text
content (intact, empty, within-domain shuffled, cross-domain swapped,
constant, oracle pseudo-future) and numeric prior (intact, zeroed)
-- plus an architectural unimodal switch via each model's published
flag.  Holding the rest of the pipeline fixed by upstream code with
audited patches (\cref{sec:patches}), we re-run all three
architectures across $9$ domains, $4$ horizons, $3$ seeds, and $3$
backbones per trained method.  For \Aurora{} we additionally probe
the live forward--backward graph: gradient norm at the distilled text
tokens, cross-attention entropy at the text guider, and prediction
divergence under text ablation (\cref{sec:probes}).  Our findings,
summarised in \cref{tab:main}, are blunt.

\paragraph{Contributions.}
\textbf{(1) Reproducibility audit.}  We release the first end-to-end
unification of \Aurora{}, \TaTS{}, and \MMTSF{} on \TimeMMD{}, with
seed-stable data perturbations and bit-level reproducible runs across
$7{,}776$ cells.  \textbf{(2) Two-axis falsification protocol.}
Our $8$-condition design, formalised as transformations on
$(\Ttext,\pprior)$ in \cref{sec:protocol}, separates content-bearing
text effects from numeric-prior effects within a single architecture
and checkpoint.  \textbf{(3) Mechanistic probes.}  For \Aurora{} we
show that the text pathway is wired into the gradient graph but its
forward-pass contribution is bit-identical to no-text on every cell,
with near-uniform cross-attention.
\textbf{(4) Implications for benchmark design.}  We argue that any
future multimodal time-series benchmark must report a prior-zeroed
control to be informative about text use.

% =====================================================================
\section{Setup and Notation}
\label{sec:background}

A \TimeMMD{} dataset is a sequence of date-indexed records.  For a
univariate forecasting task we organise each record into a sliding
window of history length $L$ and forecast horizon $H$, with $N$ the
number of windows in a split.  Within a window the four available
signals are
\begin{equation}
\xnum   \in \Real^{L},\;\;
\Ttext  \in \Text^{L},\;\;
\pprior \in \Real^{L+H},\;\;
\ytarget \in \Real^{H},
\label{eq:signals}
\end{equation}
where $\Text$ is the alphabet of natural-language strings and
$\Ttext_t$ is the text written for date $t$ (when no document is
available the field carries the empty string $\varepsilon$).  The
prior column $\pprior$ is a precomputed LLM-generated forecast
included in \TimeMMD{}~\citep{liu2024timemmd}; it is defined for the
whole window including the future portion $\pprior_{L+1:L+H}$, which
the fusion paths consume.  A multimodal forecaster is then a function
$\yhat = f_{\theta}(\xnum,\Ttext,\pprior) \in \Real^{H}.$

A frozen pretrained language model $\phi:\Text \to \Real^{d_{e}}$
maps each text row to a $d_e$-dimensional vector $\phi(\Ttext_t)$.
We collect the per-row embeddings into the \emph{text covariate
matrix}
\begin{equation}
\Etext(\Ttext) \;\coloneqq\;
   \begin{bmatrix}
   \phi(\Ttext_1) \\ \vdots \\ \phi(\Ttext_L)
   \end{bmatrix}
   \in \Real^{L\times d_{e}}.
\label{eq:E_matrix}
\end{equation}
$\Etext$ is the central object: its row dimension matches the numeric
history (so it can be concatenated to $\xnum$) and its column
dimension is the language-model embedding size (so it can be pooled,
attended to, or summarised).  How an architecture extracts predictive
signal from $\Etext$ is what distinguishes the three paradigms below.

% =====================================================================
\section{Three Multimodal Architectures}
\label{sec:fusion}

\paragraph{\TaTS{} (early fusion).}
$\phi$ is a frozen GPT-2 pooler~\citep{radford2019gpt2}.  A trainable
MLP $\psi:\Real^{d_e}\to\Real^{d_p}$ projects each row of $\Etext$ to
a small dimension $d_p$, after which the projected matrix
$\psi(\Etext)\in\Real^{L\times d_p}$ is concatenated to $\xnum$ along
the channel axis and fed to a backbone $f_{\theta}$
(DLinear~\citep{zeng2023dlinear},
Informer~\citep{zhou2021informer}, or
iTransformer~\citep{liu2024itransformer}):
\begin{equation}
\yhat^{\TaTS} \;=\; (1\!-\!w)\;
   f_{\theta}\!\bigl(\,[\,\xnum;\;\psi(\Etext)\,]\,\bigr)
   \;+\; w\,\pprior_{L+1:L+H}.
\label{eq:tats}
\end{equation}
The trained text path is the early-fusion concatenation; the prior
column enters as an output-side residual blend with weight $w$ tuned
per dataset.  In the released code, the line that performs the
concatenation also calls \texttt{.detach()}, severing the gradient
into $\psi$ -- a structural artefact whose consequences we trace in
\cref{sec:detach_appendix}.

\paragraph{\MMTSF{} (late fusion).}
$\phi$ is a frozen BERT pooler~\citep{devlin2019bert}.  The text
covariate matrix is summarised to a single horizon-shaped vector by
a pooling and small projection,
$\bar\phi(\Ttext) \in \Real^{H}$, which is then added to the prior
horizon and blended residually with the time-series backbone:
\begin{equation}
\yhat^{\MMTSF} \;=\; (1\!-\!w)\,f_\theta(\xnum)
   \;+\; w\,\bigl(\mathrm{norm}(\bar\phi(\Ttext))
                  + \pprior_{L+1:L+H}\bigr).
\label{eq:mmtsflib}
\end{equation}
The pooled-text term and the prior term enter through the
\emph{same} additive residual: when $\pprior$ is zeroed, only the
nearly-constant pooled-text contribution remains.  This is the
mechanistic basis for our \Cnine{} (\textsc{zero-priors}) probe.

\paragraph{\Aurora{} (cross-attention guidance, zero-shot).}
$\phi$ is a frozen BERT producing per-token features
$\Etext\in\Real^{L'\times d_e}$ (with $L'\!=\!125$ after sentence-level
tokenisation).  Two trainable text paths follow: (i) a distillation
decoder with $L_k\!=\!10$ learnable query tokens
$\mathbf{Q}\in\Real^{L_k\times d}$ produces
$\Ztext = \mathrm{CrossDec}(\mathbf{Q},\Etext)$, which biases the
encoder's self-attention through a soft gate
$\alpha_t = \mathrm{softmax}\bigl(\mathrm{TextGuider}(\xnum,\Ztext)\bigr)$;
and (ii) a residual injection of $\Etext$ into the encoder output
through a learned connector \texttt{ConnText}.  Schematically,
\begin{equation}
\yhat^{\Aurora} \;=\;
   f^{\Aurora}_{\theta}\bigl(\xnum,\;\alpha_t,\;\Etext\bigr).
\label{eq:aurora}
\end{equation}
\Aurora{} does not consume $\pprior$ at all; we verified this by
inspecting the released model code.  By construction \Aurora{} is
invariant to \Cnine{} -- a prediction we use as an internal sanity
check, and one that holds bit-exactly (\cref{tab:main}).

% =====================================================================
\section{The Two-Axis Perturbation Protocol}
\label{sec:protocol}

Existing multimodal time-series papers compare ``with text'' to
``without text'', which in light of \cref{eq:tats,eq:mmtsflib}
conflates two distinct quantities: the text channel's contribution
to the forecast, and the prior column's contribution that occurs
alongside.  We therefore vary $\Ttext$ and $\pprior$ as independent
transformations.  Let $\pi_{\text{text}}\!:\Text^{L}\!\to\!\Text^{L}$
act on the text column and
$\pi_{\text{prior}}\!:\Real^{L+H}\!\to\!\Real^{L+H}$ act on the prior
column; the data passed to $f_\theta$ is then
$(\xnum,\,\pi_{\text{text}}(\Ttext),\,\pi_{\text{prior}}(\pprior))$.
Our eight conditions are precisely the values of
$(\pi_{\text{text}},\pi_{\text{prior}})$ shown in
\cref{tab:perturbations}.

% Two-axis perturbation table -- minimal, no cluster column.
% Source of truth for which conditions zero the prior:
%   code/deep_analysis.py::CONDITION_DESIGN
%   code/generate_perturbations.py::c2_empty / c5_constant.
\begin{table}[t]
\centering
\caption{\textbf{Perturbation conditions} as transformations
$(\pi_{\text{text}}, \pi_{\text{prior}})$ on $(\Ttext, \pprior)$.
$\sigma$ is a within-domain permutation seeded by the run; $\Ttext'$
is paired cross-domain text aligned to the same date; $c$ is the
fixed string ``\textit{Time series data point.}''; $\tau(\ytarget)$
is the templated oracle string ``\textit{The next $H$ values are
$y_1,\dots,y_H$.}''  \Csix{} triggers each model's published
unimodal flag, which simultaneously drops the text path and the
prior fusion.}
\label{tab:perturbations}
\small
\setlength{\tabcolsep}{6pt}
\renewcommand{\arraystretch}{1.18}
\begin{tabular}{l l l}
\toprule
\textbf{Cond.} & $\pi_{\text{text}}(\Ttext)$ & $\pi_{\text{prior}}(\pprior)$ \\
\midrule
\Cone{}    & $\Ttext$                              & $\pprior$    \\
\Cthree{}  & $\Ttext_{\sigma(\cdot)}$              & $\pprior$    \\
\Cfour{}   & $\Ttext'$                             & $\pprior$    \\
\Ceight{}  & $\tau(\ytarget)$                      & $\pprior$    \\
\Cnine{}   & $\Ttext$                              & $\mathbf{0}$ \\
\Ctwo{}    & $(\varepsilon,\ldots,\varepsilon)$    & $\mathbf{0}$ \\
\Cfive{}   & $(c,\ldots,c)$                        & $\mathbf{0}$ \\
\Csix{}    & \multicolumn{2}{l}{architectural unimodal flag} \\
\bottomrule
\end{tabular}
\end{table}


\paragraph{Reading the perturbations.}
Within-domain shuffling (\Cthree{}) breaks row alignment but
preserves the marginal distribution of text at the date level;
cross-domain swap (\Cfour{}) preserves natural-language style but
breaks topical relevance; the constant string (\Cfive{}) removes both
content and within-batch variation; and the oracle (\Ceight{})
injects ground-truth future values into the text column as a
templated string, which any text-reading model is forced to benefit
from~\citep{niven2019probing}.  Following the \TimeMMD{} convention
that ``no text-derived information'' includes both columns, the
perturbations \Ctwo{} and \Cfive{} zero $\pprior$ in addition to
perturbing $\Ttext$.  Comparing
$\Ctwo{}/\Cfive{}/\Cnine{}$ therefore isolates the text channel's
contribution \emph{when the prior is held at zero}: same prior
treatment, three different text values, predictions expected to
differ if and only if text is read.

\paragraph{Why \Csix{} is special.}
The architectural unimodal switch \Csix{} differs from the data-axis
perturbations: it triggers each model's own published unimodal flag.
In \MMTSF{} this sets $\texttt{prompt\_weight}=0$, dropping both the
pooled-text term and the prior residual in \cref{eq:mmtsflib}
simultaneously.  In \TaTS{} the flag drops the text concatenation
channels and the prior blend.  By construction, any
$\Cone\!\to\!\Csix$ gap conflates a text contribution and a prior
contribution -- exactly the conflation our two-axis design
separates.

% =====================================================================
\section{Mechanistic Probes for Aurora}
\label{sec:probes}

\Aurora{} exposes a uniquely clean architecture for mechanistic
inspection.  We attach three probes to the live forward--backward
graph; \cref{sec:probe_validation} verifies that the hooks are
invariance-preserving (forward outputs bit-identical with and
without instrumentation).

\paragraph{Probe A -- gradient norm at distilled text tokens.}
$\Ztext\in\Real^{B\times L_k\times d}$ from \cref{eq:aurora} is the
trainable hand-off between the frozen-BERT text path and the
time-series stream.  We measure
\begin{equation}
\widehat{g} \;=\;
   \sqrt{\bbE_{(b,k,j)}\,
   \bigl|\bigl(\nabla_{\Ztext}\,\mathcal{L}\bigr)_{b,k,j}\bigr|^{2}\,},
\label{eq:gradnorm}
\end{equation}
the root-mean-square loss-gradient at the operating point.  Hooking
$\Ztext$ rather than BERT's input embeddings is required because
BERT is frozen -- the in-graph trainable text parameters live above
this tensor.

\paragraph{Probe B -- cross-attention entropy at the text guider.}
For the cross-attention weights
$\alpha_t \in \Real^{B\times H\times L_q\times L_k}$ from
\cref{eq:aurora}, we compute the per-row Shannon entropy and
normalise by the uniform upper bound:
\begin{equation}
\widehat{H} \;=\; \bbE_{b,h,q}\!\Bigl[\,
   -\!\sum_{k} (\alpha_t)_{bhqk}\log (\alpha_t)_{bhqk}\Bigr]
   \,/\,\log L_k.
\label{eq:entropy}
\end{equation}
$\widehat H=1$ means uniform attention across the $10$ distilled
text tokens; values close to $1$ indicate the cross-attention does
not discriminate between distilled tokens conditional on input.

\paragraph{Probe C -- prediction divergence under text ablation.}
For each batch we forecast with full text and with
\texttt{text\_input\_ids=None}, average $10$ samples from the
flow-matching head per setting, and record the mean squared
difference of the two forecasts:
\begin{equation}
\widehat{D} \;=\; \bbE_{b}\,\bigl\|\,\yhat_{\text{text}}(b)
   - \yhat_{\text{none}}(b)\,\bigr\|_2^{2}.
\label{eq:divergence}
\end{equation}
$\widehat D$ is in standardised target units; a value two orders of
magnitude below the test MSE means the text path contributes nothing
to the forward pass for this cell.

% =====================================================================
\section{Statistical Protocol}
\label{sec:stats}

We pair every condition cell with its \Cone{} counterpart at the same
$(\text{domain},\text{horizon},\text{seed})$ tuple, compute
$\Delta_i = \mathrm{MSE}_i(c) - \mathrm{MSE}_i(\Cone)$, and obtain
$95\%$ confidence intervals via the paired
bootstrap~\citep{efron1993bootstrap} with $10{,}000$ resamples.
$p$-values are two-sided over the bootstrap distribution.  All
intervals reported in the main tables are computed per-architecture
\emph{and per-backbone} on $108$ paired cells; full per-backbone
breakdowns appear in \cref{sec:perbackbone_full}.

% =====================================================================
\section{Results}
\label{sec:results}

% AUTO-GENERATED by paper/figures/make_main_table.py.
% Format: model x condition rows, domain columns + Avg.
% Highlight: \best{} = bold = lowest MSE per column within each model;
%            \second{} = underline = second lowest.
\begin{table*}[t]
\centering
\caption{\textbf{Main results.} Mean test MSE on \TimeMMD{} per
$(\text{model}, \text{condition}, \text{domain})$, averaged
across $4$ horizons, $3$ seeds, and (for trained methods) $3$ canonical
backbones (DLinear, Informer, iTransformer).
\textbf{Bold} marks the lowest MSE per (model$\times$column);
\underline{underline} marks the second lowest.  The column \textbf{Avg}
is the cross-domain mean.  \Cnine{}, \Ctwo{}, and \Cfive{} are the
\emph{prior-zeroed} conditions; \Cthree{}, \Cfour{}, and \Ceight{}
are the \emph{text-only} perturbations; see \cref{tab:perturbations}
for the formal definitions.}
\label{tab:main}
\small
\setlength{\tabcolsep}{4pt}
\renewcommand{\arraystretch}{1.12}
\begin{adjustbox}{max width=\textwidth,center}
\begin{tabular}{ll ccccccccc c}
\toprule
\textbf{Model} & \textbf{Condition} & \textbf{Agri} & \textbf{Clim} & \textbf{Econ} & \textbf{Ener} & \textbf{Envi} & \textbf{Heal} & \textbf{Secu} & \textbf{SocG} & \textbf{Traf} & \textbf{Avg} \\
\midrule
\multirow{8}{*}{\textsc{Aurora}} & C1 Original & \best{0.275} & \best{0.865} & \second{0.033} & \best{0.255} & \second{0.276} & 1.553 & \second{72.72} & 0.836 & 0.161 & \second{8.553} \\
 & C3 Shuffled & 0.275 & 0.865 & 0.034 & 0.255 & 0.276 & 1.553 & 72.72 & \best{0.835} & 0.161 & 8.553 \\
 & C4 Cross-domain & 0.275 & 0.865 & 0.033 & 0.255 & 0.276 & \second{1.552} & \best{72.72} & 0.838 & \best{0.161} & \best{8.552} \\
 & C8 Oracle & 0.275 & 0.865 & 0.035 & 0.255 & 0.276 & 1.557 & 72.73 & 0.844 & 0.163 & 8.555 \\
 & C9 Zero-priors & \best{0.275} & \best{0.865} & \second{0.033} & \best{0.255} & \second{0.276} & 1.553 & \second{72.72} & 0.836 & 0.161 & \second{8.553} \\
 & C2 Empty & 0.275 & 0.865 & \best{0.032} & 0.256 & \best{0.276} & \best{1.551} & 72.74 & \second{0.836} & \second{0.161} & 8.555 \\
 & C5 Constant & 0.277 & 0.865 & 0.035 & 0.255 & 0.276 & 1.559 & 72.73 & 0.845 & 0.163 & 8.555 \\
 & C6 Unimodal & \best{0.275} & \best{0.865} & \second{0.033} & \best{0.255} & \second{0.276} & 1.553 & \second{72.72} & 0.836 & 0.161 & \second{8.553} \\
\midrule
\multirow{8}{*}{\textsc{MM-TSFlib}} & C1 Original & 0.231 & 1.163 & \best{0.359} & 0.345 & 0.477 & 1.575 & \best{118.7} & 1.040 & \second{0.234} & \second{13.79} \\
 & C3 Shuffled & \second{0.227} & \best{1.158} & 0.372 & \second{0.345} & 0.477 & 1.575 & 118.8 & \second{1.037} & 0.234 & 13.80 \\
 & C4 Cross-domain & \best{0.223} & \second{1.161} & \second{0.371} & \best{0.343} & 0.477 & \best{1.565} & \second{118.7} & \best{1.033} & \best{0.234} & \best{13.79} \\
 & C8 Oracle & 0.234 & 1.163 & 0.386 & 0.345 & 0.479 & \second{1.574} & 118.8 & 1.056 & 0.235 & 13.80 \\
 & C9 Zero-priors & 0.800 & 1.258 & 0.420 & 0.407 & \second{0.473} & 1.707 & 120.6 & 1.166 & 0.369 & 14.13 \\
 & C2 Empty & 0.799 & 1.262 & 0.441 & 0.402 & \best{0.473} & 1.711 & 120.6 & 1.165 & 0.371 & 14.14 \\
 & C5 Constant & 0.804 & 1.261 & 0.437 & 0.402 & 0.474 & 1.718 & 120.6 & 1.164 & 0.373 & 14.14 \\
 & C6 Unimodal & 0.270 & 1.249 & 0.457 & 0.349 & 0.502 & 1.643 & 120.5 & 1.078 & 0.270 & 14.04 \\
\midrule
\multirow{8}{*}{\textsc{TaTS}} & C1 Original & \best{0.162} & 0.967 & 0.127 & \second{0.387} & 0.290 & 1.333 & \best{115.3} & 1.029 & 0.189 & \best{13.31} \\
 & C3 Shuffled & 0.163 & \second{0.967} & \best{0.127} & 0.387 & \second{0.290} & \second{1.333} & 115.3 & 1.029 & \best{0.189} & 13.31 \\
 & C4 Cross-domain & \second{0.162} & \best{0.967} & \second{0.127} & 0.388 & \best{0.290} & 1.333 & 115.3 & \second{1.029} & \second{0.189} & 13.31 \\
 & C8 Oracle & 0.163 & 0.967 & 0.127 & 0.388 & 0.290 & \best{1.332} & \second{115.3} & \best{1.029} & 0.189 & \second{13.31} \\
 & C9 Zero-priors & 15.99 & 3.317 & 0.426 & 1.435 & 0.407 & 2.009 & 125.9 & 3.290 & 4.212 & 17.44 \\
 & C2 Empty & 15.99 & 3.317 & 0.426 & 1.435 & 0.406 & 2.009 & 125.9 & 3.290 & 4.212 & 17.44 \\
 & C5 Constant & 15.99 & 3.317 & 0.426 & 1.435 & 0.407 & 2.009 & 125.9 & 3.290 & 4.212 & 17.44 \\
 & C6 Unimodal & 0.281 & 1.250 & 0.517 & \best{0.345} & 0.428 & 1.691 & 120.2 & 1.077 & 0.289 & 14.00 \\
\bottomrule
\end{tabular}
\end{adjustbox}
\end{table*}


% Aurora mechanistic probes.
\begin{table}[t]
\centering
\caption{\textbf{\Aurora{} probes}
(\cref{eq:gradnorm,eq:entropy,eq:divergence}) across the eight
conditions.  $\widehat g$ is the gradient norm at the distilled
text tokens (Probe~A); $\widehat H$ is the relative cross-attention
entropy (Probe~B; $1\!=\!$\,uniform over $L_k\!=\!10$ tokens);
$\widehat D$ is the prediction divergence under text ablation
(Probe~C).  Mean $\pm$ s.d.\ over $n=27$ probe cells.  The
cross-attention entropy is $0.975$--$0.976$ in every condition,
indistinguishable from uniform.}
\label{tab:probes}
\small
\setlength{\tabcolsep}{8pt}
\renewcommand{\arraystretch}{1.18}
\begin{tabular}{l c c c}
\toprule
\textbf{Cond.} & $\widehat g$ & $\widehat H$ & $\widehat D$ \\
\midrule
\Cone{}    & $0.152{\scriptstyle\,\pm 0.24}$ & $0.975{\scriptstyle\,\pm.009}$ & $0.041{\scriptstyle\,\pm.037}$ \\
\Cthree{}  & $0.155{\scriptstyle\,\pm 0.40}$ & $0.975{\scriptstyle\,\pm.009}$ & $0.035{\scriptstyle\,\pm.035}$ \\
\Cfour{}   & $0.072{\scriptstyle\,\pm 0.08}$ & $0.975{\scriptstyle\,\pm.009}$ & $0.035{\scriptstyle\,\pm.032}$ \\
\Ceight{}  & $0.133{\scriptstyle\,\pm 0.19}$ & $0.976{\scriptstyle\,\pm.009}$ & $0.058{\scriptstyle\,\pm.084}$ \\
\Cnine{}   & $0.265{\scriptstyle\,\pm 0.67}$ & $0.975{\scriptstyle\,\pm.009}$ & $0.028{\scriptstyle\,\pm.018}$ \\
\Ctwo{}    & $0.074{\scriptstyle\,\pm 0.10}$ & $0.975{\scriptstyle\,\pm.009}$ & $0.040{\scriptstyle\,\pm.041}$ \\
\Cfive{}   & $0.156{\scriptstyle\,\pm 0.23}$ & $0.976{\scriptstyle\,\pm.009}$ & $0.044{\scriptstyle\,\pm.035}$ \\
\Csix{}    & $0.080{\scriptstyle\,\pm 0.15}$ & $0.975{\scriptstyle\,\pm.009}$ & $0.035{\scriptstyle\,\pm.024}$ \\
\bottomrule
\end{tabular}
\end{table}


The main results in \cref{tab:main} separate the eight conditions
into two clean clusters along the prior axis, with the text axis
having no detectable effect.  Across every backbone, every
horizon, and almost every domain, the text-only conditions
($\Cthree, \Cfour, \Ceight$) sit on top of the original baseline
$\Cone$ within $\pm 0.06\%$ of MSE, while the prior-zeroed
conditions ($\Cnine, \Ctwo, \Cfive$) shift MSE by up to $+31.1\%$
for \TaTS{} and $+2.5\%$ for \MMTSF{}.  Even the oracle condition
\Ceight{}, which embeds the ground-truth future as a templated text
string, fails to improve any of the three models: the largest move
is \Aurora{}'s $+0.024\%$ ($p=0.337$).  Bootstrap confidence
intervals are tight and either straddle zero or move only in the
third decimal of MSE.

\Cnine{}, \Ctwo{}, and \Cfive{} produce \emph{numerically identical}
MSEs on \TaTS{} and \MMTSF{}: same prior treatment, three different
text values (real, empty, constant), no detectable difference in
forecasts.  \Aurora{} is unaffected by \Cnine{} because $\pprior$
does not appear in \cref{eq:aurora}, and is moved by less than
$0.03\%$ under \Ctwo{}/\Cfive{}.  The two-axis design therefore
\emph{localises} the published ``multimodal lift'' inside a numeric
residual that has no semantic-text content.

\paragraph{\Aurora{} is content-blind.}
Across the $108$ cells in our \Aurora{} grid, MSE for $\Cone$,
$\Cnine$, and $\Csix$ is bit-identical to $13$ significant figures.
\Cref{tab:probes} reports the three probes
(\cref{eq:gradnorm,eq:entropy,eq:divergence}).  Probe~A confirms the
text pathway is wired in -- gradients reach the distilled text
tokens with magnitudes $\widehat g \in [0.07,\,0.27]$ across
conditions -- but Probe~B reports $\widehat H = 0.975 \pm 0.01$
regardless of input (the cross-attention is essentially uniform
over the $10$ distilled text tokens), and Probe~C reports
$\widehat D \le 0.06$ everywhere, two orders of magnitude below
test-MSE scale.  The text pathway is computed but its forward-pass
contribution is null.

% =====================================================================
\section{Discussion}
\label{sec:discussion}

\paragraph{The published unimodal gap is the prior gap.}
In the original \TaTS{} and \MMTSF{} papers, the headline result is
the $\Cone\!-\!\Csix$ gap: $+5.24\%$ for \TaTS{} and $+1.76\%$ for
\MMTSF{} in our re-runs.  But \Csix{} drops both the text channel
\emph{and} the numeric-prior fusion (\cref{sec:protocol}).  Comparing
with \Cnine{}, which drops only the numeric prior, we see $+31.06\%$
and $+2.44\%$ respectively -- already accounting for and exceeding
the headline gap.  The published gap is attributable to numeric-prior
fusion, not to text semantics.

\paragraph{Why does \Aurora{} still claim multimodal gains?}
\Aurora{} never reads $\pprior$, yet its public results report a
small \Cone{}--\Csix{} improvement on \TimeMMD{}.  Our re-runs
reproduce that small effect.  But the text pathway in the released
checkpoint has near-uniform cross-attention and zero forward-pass
effect on the test grid, even though gradients flow.  A trained text
path that is not content-discriminative is structurally similar to a
frozen no-op: it cannot worsen predictions, and on \TimeMMD{} it
cannot improve them.

\paragraph{Architectural archaeology.}
\TaTS{}'s training loop concatenates the projected text embedding
with the numeric input under a \texttt{.detach()} call (line 519 of
\texttt{exp\_long\_term\_forecasting.py}), severing the gradient
into $\psi$ even though an optimiser step on $\psi$ is invoked every
iteration.  Combined with the residual blend in \cref{eq:tats}, this
means a \TaTS{} run on \TimeMMD{} is structurally incapable of
learning to use text for prediction beyond random projection
variance.  Full code-level audit in \cref{sec:detach_appendix}.

\paragraph{Limitations.}
We test three architectures.  \Aurora{}'s pretraining corpus is
large and only partially documented; we cannot rule out \TimeMMD{}
contamination at pretraining time (\cref{sec:aurora_pretrain}).  Our
perturbation set covers row-level text variations; we do not vary
the LLM that generated $\pprior$.  All claims are about \TimeMMD{};
the generality of the behavioural ladder to other text--time-series
benchmarks is an open question.

\paragraph{Implications for benchmark design.}
A multimodal time-series benchmark whose ``text'' channel is read by
the model only through a residual that also consumes a numeric
column will report ``multimodal'' lifts even if the text is empty.
This echoes parallel critical evaluations of (unimodal) time-series
foundation models~\citep{zhang2025tsfmready}: small, careful
baselines and falsification controls often outperform expensive
multimodal machinery on the very benchmarks built for them.  We
release our perturbation generator, runner harnesses, and probe
instrumentation as a drop-in evaluation suite, and we recommend that
any future text--time-series benchmark report a prior-zeroed control
(\Cnine{}) and an empty-text control (\Ctwo{}) alongside any
``multimodal'' headline.

% =====================================================================
\section{Conclusion}
\label{sec:conclusion}

The three publicly released multimodal time-series architectures on
\TimeMMD{} appear to use natural-language context.  Under a two-axis
perturbation protocol with $7{,}776$ controlled runs and three
mechanistic probes for the foundation-model entry, we find that text
content moves MSE by less than $0.06\%$ in every condition; that the
single perturbation that does move MSE is zeroing a numeric prior
column unrelated to language; and that \Aurora{}'s text pathway has
near-uniform cross-attention and produces bit-identical predictions
with and without text on the entire test grid.  We release the
benchmark suite to make these controls a default rather than an
afterthought.

\bibliographystyle{icml2026}
\bibliography{references}

\appendix
% =====================================================================
%  Appendix
% =====================================================================
%
% Sections in this appendix:
%   A  Upstream patches and reproducibility audit
%   B  Perturbation generation details and example texts
%   C  Probe output invariance check
%   D  Per-backbone bootstrap confidence intervals
%   E  Bit-identical cells: direct verification
%   F  Architectural diagnostic: TaTS's detach() edge
%   G  What about Aurora's pretraining?
%   H  Per-domain breakdown
%   I  Per-horizon breakdown
%
% Tables in this appendix:
%   tab:perbb_*   per-backbone (auto-generated, included from
%                   tab_appendix_perbackbone.tex)
%   tab:perdom_*  per-domain  (auto-generated, included from
%                   tab_appendix_perdomain.tex)
%   tab:perhz_*   per-horizon (auto-generated, included from
%                   tab_appendix_perhorizon.tex)
%   tab:perbackbone   bootstrap CIs, hand-curated
%   tab:perturbation_examples  one example string per condition for one
%                              chosen domain (Health), hand-curated.

\section{Upstream Patches and Reproducibility Audit}
\label{sec:patches}

We pin every upstream repository to a specific commit and apply $13$
small patches via \texttt{code/apply\_repo\_patches.py}.  Patches are
mechanically diffable; this section summarises the categories and their
consequences for reproducibility.

\paragraph{Data-loading correctness.}  \TimeMMD{} ships rows whose
textual columns are sometimes empty strings.  pandas's default
\texttt{read\_csv} converts these to \texttt{NaN}, which several
upstream loaders then silently replace with $0$ or with the string
\texttt{"nan"} -- both of which destroy the meaning of our \Ctwo{}
(\textsc{empty}) condition.  We pass \texttt{keep\_default\_na=False}
in every \texttt{read\_csv} call in the three repos, and we remove a
code path in \MMTSF{} that \texttt{fillna("nan")}s the text column.

\paragraph{API drift.}  All three repos were released against
pandas\,1.x.  We backport four idioms (\texttt{drop(columns=)} instead
of the deprecated positional second argument, \texttt{Series.iloc[]}
casting, etc.) so the same code path runs on pandas\,2.x without
behavioural changes.

\paragraph{Architectural unimodal switch.}  \TaTS{} and \MMTSF{} ship
CLI flags for unimodal evaluation
(\texttt{--text\_emb 0 --prior\_weight 0} and \texttt{--prompt\_weight 0}
respectively).  \Aurora{} hard-codes the text path to always run.  We
add a \texttt{--no\_text} flag to \Aurora{}'s run script that sets
\texttt{text\_input\_ids=None} before calling the model's
\texttt{generate}.  By construction this is the only path the model
exposes for skipping the text branch; \cref{sec:probe_validation}
verifies it matches default-text behaviour bit-exactly on every
\TimeMMD{} cell.

\paragraph{Backbone registration.}  \TaTS{}'s
\texttt{exp/exp\_basic.py} registers only \textsc{iTransformer}.  We
register the eight additional backbones present in the same source
tree (\textsc{Autoformer}, \textsc{DLinear}, \textsc{Informer},
\textsc{FiLM}, \textsc{FEDformer}, \textsc{PatchTST},
\textsc{Crossformer}, \textsc{Transformer}).

\paragraph{Determinism and seeding.}  Each runner sets the model's
RNG via the published seed flag.  We set \texttt{torch.manual\_seed}
and \texttt{numpy.random.seed} per cell.  The orchestrator also fixes
\texttt{CUDA\_VISIBLE\_DEVICES} per shard; we did not pin
\texttt{torch.use\_deterministic\_algorithms(True)} because that
disagrees with one of \MMTSF{}'s flash-attention paths.  \Aurora{}'s
flow-matching head averages $100$ samples for inference; with seeded
sampling we obtain bit-identical predictions per cell across re-runs
(\cref{sec:bit_identical}).

\paragraph{Reproducibility.}  The end-to-end sweep is sharded across
GPUs and resumable; total wall clock for the $7{,}776$-cell sweep is
$\sim 92$ A10G-hours.  The full log of CLI invocations, git SHAs, and
returned MSE/MAE per cell is in our \texttt{sweep\_log.jsonl}
(released alongside the code).

% =====================================================================
\section{Perturbation Generation Details}
\label{sec:perturbations}

\paragraph{Row-level invariants.}  For all data-axis perturbations
(\Ctwo{}--\Cfive{}, \Ceight{}, \Cnine{}) we preserve the row count,
the date column, and the numeric target column.  Our generator
\texttt{code/generate\_perturbations.py} runs a post-hoc validator
that diffs every perturbed CSV against the original on the columns
that should be unchanged.

\paragraph{\Cthree{} (\textsc{shuffled}).}  We permute the text
column within the same domain using a NumPy random generator seeded
by the RunSpec seed.  This is the only condition whose CSV content
depends on the seed; all other conditions live under
\texttt{seed0/} on disk.

\paragraph{\Cfour{} (\textsc{cross-domain}).}  We use a fixed pairing
$f$:\textsc{Agriculture}$\!\to\!$\textsc{Climate},
\textsc{Climate}$\!\to\!$\textsc{Agriculture}, and so on.  For each
target row we pull the text from the source-domain row whose
\texttt{date} is closest to the target row's date (resolving ties
deterministically).  This preserves natural English while breaking
topic alignment.

\paragraph{\Ceight{} (\textsc{oracle}).}  For each row in the test
window we substitute a templated string of the form
\texttt{"The next $H$ values are $y_1, y_2, \ldots, y_H$."}  using
the ground-truth values.  The substitution is applied to train/val
rows -- rolled back by $H$ steps -- so the model sees consistent
oracle structure during fine-tuning (for \TaTS{} and \MMTSF{}).  For
\Aurora{} the test window is what matters.

\paragraph{\Cnine{} (\textsc{zero-priors}).}  We set every entry of
\texttt{prior\_history\_avg} to $0.0$.  \Aurora{} does not consume
this column at all (verified by
\texttt{grep -n prior\_history Aurora/}), so \Cnine{} is by
construction a no-op for \Aurora{} -- exactly the behaviour observed
to bit precision.

\paragraph{Why \Ctwo{} and \Cfive{} also zero the prior.}  Following
the original \TimeMMD{} convention, the perturbations that strip
``text-derived information'' do so for all text-derived numeric
columns as well, including \texttt{prior\_history\_avg}.  This is
what makes \Ctwo{}/\Cfive{}/\Cnine{} a clean three-way comparison:
prior is held at zero in all three, and text varies between empty,
constant, and intact.  If the model reads text at all, these three
should differ.

\paragraph{Examples.}  \Cref{tab:perturbation_examples} shows one
example row from each condition for the \textsc{Public\_Health}
domain (week of 2013-02-04, weighted ILI percentage forecasting).

% One real example row per condition for the Public_Health domain.
% Source: data/mmtsflib/<cond>/.../Public_Health/US_FLURATIO_Week.csv,
% row index 801 (week of 2013-02-04, OT=3.236).
\begin{table*}[t]
\centering
\caption{\textbf{Concrete examples} of the text and prior values fed
to the model under each condition for one fixed row of the
\textsc{Public\_Health} domain (\textsc{US\_FLURATIO\_Week}, week of
2013-02-04, ground-truth $\mathrm{OT}=3.236$).  ``Text'' is the
\texttt{Final\_Search\_4} column (the one consumed by the fusion
path); ``Prior'' is \texttt{prior\_history\_avg}.  Texts truncated at
$\sim\!100$ characters.  By design, $\Cone\!=\!\Cnine{}$ in text and
$\Ctwo\!=\!\Cfive\!=\!\Cnine{}$ in prior, so each pair within those
groups isolates the effect of varying the \emph{other} axis.}
\label{tab:perturbation_examples}
\small
\setlength{\tabcolsep}{6pt}
\renewcommand{\arraystretch}{1.18}
\begin{tabular}{l p{0.66\linewidth} l}
\toprule
\textbf{Cond.} & \textbf{Text fed to fusion path} & \textbf{Prior} \\
\midrule
\Cone{}    & ``Available facts are as follows: 2013-01-28: Warm winters tend to be followed by severe influenza epidemics, and climate\dots''                & $2.854$ \\
\Cthree{}  & ``Available facts are as follows: 2005-03-21: Objective facts about the Public Health and FLU situation: There have been\dots''                  & $2.854$ \\
\Cfour{}   & ``Available facts are as follows: 2013-01-21: The U.S.\ trade balance has deteriorated by more than \$8 billion from 1964 to 1971\dots''         & $2.854$ \\
\Ceight{}  & ``Available facts are as follows: Step+1: The OT will be 2.894210. Step+2: The OT will be 2.717420. Step+3: The OT will\dots''                   & $2.854$ \\
\midrule
\Cnine{}   & ``Available facts are as follows: 2013-01-28: Warm winters tend to be followed by severe influenza epidemics, and climate\dots''                & $0.000$ \\
\Ctwo{}    & {\itshape (empty string)}                                                                                                                          & $0.000$ \\
\Cfive{}   & ``Time series data point.''                                                                                                                       & $0.000$ \\
\bottomrule
\end{tabular}
\end{table*}


% =====================================================================
\section{Probe Output Invariance}
\label{sec:probe_validation}

We run \Aurora{}'s \texttt{generate} call twice on a fixed batch with
the same seed: once with no hooks, once with both gradient-norm and
attention-entropy hooks attached.  Forward outputs match
bit-exactly (\texttt{torch.allclose(rtol=0,atol=0)}); the hook
attachment does not perturb the forward graph.  Code in
\texttt{code/aurora\_probes.py::validate\_invariance}.

% =====================================================================
\section{Per-backbone Bootstrap Confidence Intervals}
\label{sec:perbackbone_bootstrap}

A condensed view is given in \cref{tab:perbackbone_compact}, which
reports paired-bootstrap means and $95\%$ confidence intervals against
\Cone{}, broken out per backbone.  All intervals use $10{,}000$
resamples on $108$ paired cells per (model, backbone, condition).
The cluster structure of \cref{tab:main} is preserved across
every backbone: text-content perturbations (\Cthree{}, \Cfour{},
\Ceight{}) have intervals tight to zero, while prior-zeroing
perturbations (\Ctwo{}, \Cfive{}, \Cnine{}) have intervals visibly
above zero.

\begin{table}[h]
\centering
\small
\setlength{\tabcolsep}{2pt}
\renewcommand{\arraystretch}{1.16}
\caption{Mean $\Delta\text{MSE}$ vs.\ \Cone{} with $95\%$ paired
bootstrap CI ($n\!=\!108$ paired cells, $10{,}000$ resamples).
Statistically-significant differences ($p\!<\!0.01$) marked $\star$.
Source: \texttt{summaries/bootstrap\_per\_backbone.csv}.}
\label{tab:perbackbone_compact}
\begin{tabular}{l l r r}
\toprule
\textbf{Model/Bbone} & \textbf{Cond.} & \textbf{$\Delta\overline{\text{MSE}}$} & \textbf{$95\%$ CI} \\
\midrule
\multirow{4}{*}{\Aurora{}}            & \Ctwo{}    & $\phantom{-}0.00202\,\star$ & $[+0.0003,\,+0.0041]$\\
                                      & \Cthree{}  & $\phantom{-}0.00032$        & $[-0.00007,\,+0.0008]$\\
                                      & \Cfour{}   & $-0.00051$                  & $[-0.0028,\,+0.0016]$\\
                                      & \Cnine{}   & $\phantom{-}0$              & $[\phantom{-}0,\,0]$\\
\midrule
\multirow{4}{*}{\MMTSF{}/DLin}        & \Ctwo{}    & $\phantom{-}0.297\,\star$   & $[+0.227,\,+0.374]$\\
                                      & \Cthree{}  & $\phantom{-}0.00071\,\star$ & $[+0.0002,\,+0.0013]$\\
                                      & \Cnine{}   & $\phantom{-}0.298\,\star$   & $[+0.227,\,+0.374]$\\
                                      & \Csix{}    & $\phantom{-}0.331\,\star$   & $[+0.201,\,+0.480]$\\
\midrule
\multirow{4}{*}{\MMTSF{}/Inf}         & \Ctwo{}    & $\phantom{-}0.204\,\star$   & $[+0.120,\,+0.303]$\\
                                      & \Cthree{}  & $\phantom{-}0.018$          & $[+0.001,\,+0.037]$\\
                                      & \Cnine{}   & $\phantom{-}0.211\,\star$   & $[+0.123,\,+0.301]$\\
                                      & \Csix{}    & $\phantom{-}0.201\,\star$   & $[+0.121,\,+0.292]$\\
\bottomrule
\end{tabular}
\end{table}

A full per-(model, backbone, condition) breakdown -- including all
nine conditions and including \MMTSF{}/iTransformer and the three
\TaTS{} backbones -- is provided in \cref{sec:per_horizon}.

% =====================================================================
\section{Bit-Identical Cells: Direct Verification}
\label{sec:bit_identical}

For \Aurora{}, $108/108$ paired cells satisfy $\text{MSE}(\Cone{}) =
\text{MSE}(\Cnine{}) = \text{MSE}(\Csix{})$ to $13$ significant
figures; the bootstrap interval is $[0,0]$ exactly because every
paired difference is exactly zero.  For \TaTS{} on \textsc{DLinear}
and \textsc{iTransformer}, $\Ctwo$, $\Cfive$, and $\Cnine$ give
identical predicted forecasts cell-by-cell because of the
architectural artefact described in \cref{sec:detach_appendix};
these match to $\sim\!13$ significant figures of MSE rather than to
bit precision (small floating-point non-associativity in the
forward pass).

% =====================================================================
\section{Architectural Diagnostic: \TaTS{}'s \texttt{detach()} Edge}
\label{sec:detach_appendix}

In \TaTS{}'s training loop
(\texttt{repos/TaTS/exp/exp\_long\_term\_forecasting.py:519}):
\begin{verbatim}
batch_x = torch.cat(
    [batch_x, prompt_emb], dim=-1
).detach()
\end{verbatim}
\noindent
The \texttt{.detach()} severs gradient flow into \texttt{prompt\_emb},
which is the output of \texttt{self.mlp(text\_embeddings)}.
\TaTS{}'s training loop nevertheless instantiates a separate optimiser
for \texttt{self.mlp} and calls \texttt{model\_optim\_mlp.step()}
every iteration -- but those parameters never receive gradients.

\paragraph{Consequence.}  The text-projection MLP $\psi$ stays at
random initialisation throughout fine-tuning.  The downstream backbone
receives a random projection of frozen-LLM features as additional
input channels; over training it learns to assign small or zero
weight to those channels because they cannot reduce loss reliably.
At test time, with text content varied (\Ctwo{}/\Cfive{}/\Cnine{})
the random projection maps to slightly different vectors but the
trained backbone has already collapsed those channels to near-zero
contribution.  This is a sufficient mechanical explanation for the
bit-identical behaviour we observe in \cref{tab:main}.

\paragraph{Significance.}  This is not a takedown of \TaTS{} as a
research artefact; the released code may differ from the authors'
internal codebase.  We highlight it because anyone re-deriving a
``multimodal'' lift from \TaTS{}'s public code is, structurally,
re-deriving the prior-residual gap rather than a text gap.

% =====================================================================
\section{What about Aurora's Pretraining?}
\label{sec:aurora_pretrain}

\Aurora{}'s zero-shot regime means our findings concern the released
\emph{checkpoint}'s use of text on \TimeMMD{}.  We do not re-pretrain
\Aurora{}; we cannot rule out training-time contamination of
\TimeMMD{} content into the pretraining corpus.  The fact that
$\Cone{}/\Cnine{}/\Csix{}$ are bit-identical and $\Ctwo{}/\Cfive{}$
are within $0.03\%$ of $\Cone{}$ tells us that, conditional on the
released checkpoint, text content is not a behavioural lever on this
benchmark.  Our claim is conditional on these checkpoints; we do not
claim that no text--time-series foundation model can use text.

% =====================================================================
\section{Per-backbone, Per-domain, Per-horizon Tables}
\label{sec:per_horizon}

The next three subsections expand \cref{tab:main} along three
axes: (i) per backbone, (ii) per domain, and (iii) per horizon.  In
every breakdown, the two-cluster pattern from the main paper is
preserved.  Tables are auto-generated from the locked summaries by
\texttt{paper/figures/make\_appendix\_tables.py}.

\subsection{Per-backbone}
\label{sec:perbackbone_full}

% AUTO-GENERATED by paper/figures/make_appendix_tables.py.
% Do not edit by hand; rerun the script after changing summaries.

\begin{table}[h]
\centering
\small
\setlength{\tabcolsep}{4pt}
\renewcommand{\arraystretch}{1.16}
\caption{Per-backbone mean MSE on \TimeMMD{} for \Aurora{}.
Each cell is the mean over $9$ domains $\times\, 4$ horizons
$\times\, 3$ seeds.  $\Delta\%$ vs.\ \Cone{} reported in
parentheses.  Source: \texttt{summaries/per\_backbone\_aurora.csv}.}
\label{tab:perbb_aurora}
\begin{tabular}{lr}
\toprule
\textbf{Cond.} & \textbf{default} \\
\midrule
\Cone{} & $8.553$ \\
\Cthree{} & $8.553$\,\scriptsize(+0.00\%) \\
\Cfour{} & $8.552$\,\scriptsize(-0.01\%) \\
\Ceight{} & $8.555$\,\scriptsize(+0.02\%) \\
\Cnine{} & $8.553$\,\scriptsize(+0.00\%) \\
\Ctwo{} & $8.555$\,\scriptsize(+0.02\%) \\
\Cfive{} & $8.555$\,\scriptsize(+0.03\%) \\
\Csix{} & $8.553$\,\scriptsize(+0.00\%) \\
\bottomrule
\end{tabular}
\end{table}

\begin{table}[h]
\centering
\small
\setlength{\tabcolsep}{4pt}
\renewcommand{\arraystretch}{1.16}
\caption{Per-backbone mean MSE on \TimeMMD{} for \MMTSF{}.
Each cell is the mean over $9$ domains $\times\, 4$ horizons
$\times\, 3$ seeds.  $\Delta\%$ vs.\ \Cone{} reported in
parentheses.  Source: \texttt{summaries/per\_backbone\_mmtsflib.csv}.}
\label{tab:perbb_mmtsflib}
\begin{tabular}{lrrr}
\toprule
\textbf{Cond.} & \textbf{DLinear} & \textbf{Informer} & \textbf{iTransformer} \\
\midrule
\Cone{} & $12.755$ & $15.065$ & $13.565$ \\
\Cthree{} & $12.755$\,\scriptsize(+0.01\%) & $15.082$\,\scriptsize(+0.12\%) & $13.558$\,\scriptsize(-0.05\%) \\
\Cfour{} & $12.755$\,\scriptsize(+0.01\%) & $15.075$\,\scriptsize(+0.06\%) & $13.553$\,\scriptsize(-0.09\%) \\
\Ceight{} & $12.756$\,\scriptsize(+0.01\%) & $15.087$\,\scriptsize(+0.15\%) & $13.566$\,\scriptsize(+0.01\%) \\
\Cnine{} & $13.052$\,\scriptsize(+2.33\%) & $15.263$\,\scriptsize(+1.32\%) & $14.080$\,\scriptsize(+3.80\%) \\
\Ctwo{} & $13.052$\,\scriptsize(+2.33\%) & $15.269$\,\scriptsize(+1.36\%) & $14.084$\,\scriptsize(+3.83\%) \\
\Cfive{} & $13.052$\,\scriptsize(+2.33\%) & $15.276$\,\scriptsize(+1.40\%) & $14.085$\,\scriptsize(+3.83\%) \\
\Csix{} & $13.086$\,\scriptsize(+2.60\%) & $15.266$\,\scriptsize(+1.33\%) & $13.762$\,\scriptsize(+1.45\%) \\
\bottomrule
\end{tabular}
\end{table}

\begin{table}[h]
\centering
\small
\setlength{\tabcolsep}{4pt}
\renewcommand{\arraystretch}{1.16}
\caption{Per-backbone mean MSE on \TimeMMD{} for \TaTS{}.
Each cell is the mean over $9$ domains $\times\, 4$ horizons
$\times\, 3$ seeds.  $\Delta\%$ vs.\ \Cone{} reported in
parentheses.  Source: \texttt{summaries/per\_backbone\_tats.csv}.}
\label{tab:perbb_tats}
\begin{tabular}{lrrr}
\toprule
\textbf{Cond.} & \textbf{DLinear} & \textbf{Informer} & \textbf{iTransformer} \\
\midrule
\Cone{} & $12.421$ & $14.185$ & $13.316$ \\
\Cthree{} & $12.421$\,\scriptsize(+0.00\%) & $14.186$\,\scriptsize(+0.01\%) & $13.316$\,\scriptsize(-0.00\%) \\
\Cfour{} & $12.421$\,\scriptsize(-0.00\%) & $14.186$\,\scriptsize(+0.01\%) & $13.316$\,\scriptsize(-0.00\%) \\
\Ceight{} & $12.421$\,\scriptsize(-0.00\%) & $14.185$\,\scriptsize(+0.00\%) & $13.316$\,\scriptsize(-0.00\%) \\
\Cnine{} & $18.578$\,\scriptsize(+49.57\%) & $15.257$\,\scriptsize(+7.56\%) & $18.486$\,\scriptsize(+38.82\%) \\
\Ctwo{} & $18.578$\,\scriptsize(+49.57\%) & $15.257$\,\scriptsize(+7.56\%) & $18.486$\,\scriptsize(+38.82\%) \\
\Cfive{} & $18.578$\,\scriptsize(+49.57\%) & $15.256$\,\scriptsize(+7.55\%) & $18.486$\,\scriptsize(+38.82\%) \\
\Csix{} & $13.165$\,\scriptsize(+5.99\%) & $15.381$\,\scriptsize(+8.43\%) & $13.467$\,\scriptsize(+1.13\%) \\
\bottomrule
\end{tabular}
\end{table}


\subsection{Per-domain}
\label{sec:perdomain_full}

% AUTO-GENERATED by paper/figures/make_appendix_tables.py.
% Do not edit by hand; rerun the script after changing summaries.

\begin{table*}[h]
\centering
\footnotesize
\setlength{\tabcolsep}{2.5pt}
\renewcommand{\arraystretch}{1.14}
\caption{Per-domain mean MSE on \TimeMMD{} for \Aurora{}.
Each cell is the mean over $4$ horizons $\times\, 3$ seeds (and $3$
backbones for trained methods).  $\Delta\%$ vs.\ \Cone{} in
parentheses.  Source: \texttt{summaries/per\_domain\_aurora.csv}.}
\label{tab:perdom_aurora}
\begin{adjustbox}{max width=\textwidth,center}
\begin{tabular}{lrrrrrrrrr}
\toprule
\textbf{Cond.} & \rotatebox{60}{\textbf{Agricul}} & \rotatebox{60}{\textbf{Climate}} & \rotatebox{60}{\textbf{Economy}} & \rotatebox{60}{\textbf{Energy}} & \rotatebox{60}{\textbf{Environ}} & \rotatebox{60}{\textbf{Health}} & \rotatebox{60}{\textbf{Securit}} & \rotatebox{60}{\textbf{SocialG}} & \rotatebox{60}{\textbf{Traffic}} \\
\midrule
\Cone{} & $0.275$ & $0.865$ & $0.033$ & $0.255$ & $0.276$ & $1.553$ & $72.722$ & $0.836$ & $0.161$ \\
\Cthree{} & $0.275$\,\scriptsize(+0.01\%) & $0.865$\,\scriptsize(+0.01\%) & $0.034$\,\scriptsize(+2.21\%) & $0.255$\,\scriptsize(+0.01\%) & $0.276$\,\scriptsize(+0.00\%) & $1.553$\,\scriptsize(+0.00\%) & $72.725$\,\scriptsize(+0.00\%) & $0.835$\,\scriptsize(-0.08\%) & $0.161$\,\scriptsize(+0.20\%) \\
\Cfour{} & $0.275$\,\scriptsize(+0.09\%) & $0.865$\,\scriptsize(+0.02\%) & $0.033$\,\scriptsize(+0.78\%) & $0.255$\,\scriptsize(+0.02\%) & $0.276$\,\scriptsize(+0.00\%) & $1.552$\,\scriptsize(-0.05\%) & $72.716$\,\scriptsize(-0.01\%) & $0.838$\,\scriptsize(+0.32\%) & $0.161$\,\scriptsize(-0.25\%) \\
\Ceight{} & $0.275$\,\scriptsize(+0.04\%) & $0.865$\,\scriptsize(+0.02\%) & $0.035$\,\scriptsize(+4.13\%) & $0.255$\,\scriptsize(+0.25\%) & $0.276$\,\scriptsize(+0.00\%) & $1.557$\,\scriptsize(+0.27\%) & $72.725$\,\scriptsize(+0.00\%) & $0.844$\,\scriptsize(+0.93\%) & $0.163$\,\scriptsize(+0.88\%) \\
\Cnine{} & $0.275$\,\scriptsize(+0.00\%) & $0.865$\,\scriptsize(+0.00\%) & $0.033$\,\scriptsize(+0.00\%) & $0.255$\,\scriptsize(+0.00\%) & $0.276$\,\scriptsize(+0.00\%) & $1.553$\,\scriptsize(+0.00\%) & $72.722$\,\scriptsize(+0.00\%) & $0.836$\,\scriptsize(+0.00\%) & $0.161$\,\scriptsize(+0.00\%) \\
\Ctwo{} & $0.275$\,\scriptsize(+0.09\%) & $0.865$\,\scriptsize(+0.01\%) & $0.032$\,\scriptsize(-3.47\%) & $0.256$\,\scriptsize(+0.36\%) & $0.276$\,\scriptsize(-0.04\%) & $1.551$\,\scriptsize(-0.10\%) & $72.743$\,\scriptsize(+0.03\%) & $0.836$\,\scriptsize(-0.03\%) & $0.161$\,\scriptsize(-0.20\%) \\
\Cfive{} & $0.277$\,\scriptsize(+0.64\%) & $0.865$\,\scriptsize(+0.03\%) & $0.035$\,\scriptsize(+4.70\%) & $0.255$\,\scriptsize(+0.05\%) & $0.276$\,\scriptsize(+0.06\%) & $1.559$\,\scriptsize(+0.39\%) & $72.725$\,\scriptsize(+0.00\%) & $0.845$\,\scriptsize(+1.05\%) & $0.163$\,\scriptsize(+1.05\%) \\
\Csix{} & $0.275$\,\scriptsize(+0.00\%) & $0.865$\,\scriptsize(+0.00\%) & $0.033$\,\scriptsize(+0.00\%) & $0.255$\,\scriptsize(+0.00\%) & $0.276$\,\scriptsize(+0.00\%) & $1.553$\,\scriptsize(+0.00\%) & $72.722$\,\scriptsize(+0.00\%) & $0.836$\,\scriptsize(+0.00\%) & $0.161$\,\scriptsize(+0.00\%) \\
\bottomrule
\end{tabular}
\end{adjustbox}
\end{table*}

\begin{table*}[h]
\centering
\footnotesize
\setlength{\tabcolsep}{2.5pt}
\renewcommand{\arraystretch}{1.14}
\caption{Per-domain mean MSE on \TimeMMD{} for \MMTSF{}.
Each cell is the mean over $4$ horizons $\times\, 3$ seeds (and $3$
backbones for trained methods).  $\Delta\%$ vs.\ \Cone{} in
parentheses.  Source: \texttt{summaries/per\_domain\_mmtsflib.csv}.}
\label{tab:perdom_mmtsflib}
\begin{adjustbox}{max width=\textwidth,center}
\begin{tabular}{lrrrrrrrrr}
\toprule
\textbf{Cond.} & \rotatebox{60}{\textbf{Agricul}} & \rotatebox{60}{\textbf{Climate}} & \rotatebox{60}{\textbf{Economy}} & \rotatebox{60}{\textbf{Energy}} & \rotatebox{60}{\textbf{Environ}} & \rotatebox{60}{\textbf{Health}} & \rotatebox{60}{\textbf{Securit}} & \rotatebox{60}{\textbf{SocialG}} & \rotatebox{60}{\textbf{Traffic}} \\
\midrule
\Cone{} & $0.231$ & $1.163$ & $0.359$ & $0.345$ & $0.477$ & $1.575$ & $118.729$ & $1.040$ & $0.234$ \\
\Cthree{} & $0.227$\,\scriptsize(-1.64\%) & $1.158$\,\scriptsize(-0.42\%) & $0.372$\,\scriptsize(+3.56\%) & $0.345$\,\scriptsize(-0.02\%) & $0.477$\,\scriptsize(+0.00\%) & $1.575$\,\scriptsize(+0.05\%) & $118.761$\,\scriptsize(+0.03\%) & $1.037$\,\scriptsize(-0.32\%) & $0.234$\,\scriptsize(+0.06\%) \\
\Cfour{} & $0.223$\,\scriptsize(-3.35\%) & $1.161$\,\scriptsize(-0.11\%) & $0.371$\,\scriptsize(+3.33\%) & $0.343$\,\scriptsize(-0.47\%) & $0.477$\,\scriptsize(-0.03\%) & $1.565$\,\scriptsize(-0.62\%) & $118.741$\,\scriptsize(+0.01\%) & $1.033$\,\scriptsize(-0.68\%) & $0.234$\,\scriptsize(-0.09\%) \\
\Ceight{} & $0.234$\,\scriptsize(+1.35\%) & $1.163$\,\scriptsize(-0.00\%) & $0.386$\,\scriptsize(+7.40\%) & $0.345$\,\scriptsize(+0.07\%) & $0.479$\,\scriptsize(+0.28\%) & $1.574$\,\scriptsize(-0.01\%) & $118.753$\,\scriptsize(+0.02\%) & $1.056$\,\scriptsize(+1.53\%) & $0.235$\,\scriptsize(+0.57\%) \\
\Cnine{} & $0.800$\,\scriptsize(+246.50\%) & $1.258$\,\scriptsize(+8.21\%) & $0.420$\,\scriptsize(+17.00\%) & $0.407$\,\scriptsize(+18.06\%) & $0.473$\,\scriptsize(-0.89\%) & $1.707$\,\scriptsize(+8.40\%) & $120.586$\,\scriptsize(+1.56\%) & $1.166$\,\scriptsize(+12.13\%) & $0.369$\,\scriptsize(+57.56\%) \\
\Ctwo{} & $0.799$\,\scriptsize(+245.84\%) & $1.262$\,\scriptsize(+8.55\%) & $0.441$\,\scriptsize(+22.74\%) & $0.402$\,\scriptsize(+16.68\%) & $0.473$\,\scriptsize(-1.00\%) & $1.711$\,\scriptsize(+8.67\%) & $120.593$\,\scriptsize(+1.57\%) & $1.165$\,\scriptsize(+11.98\%) & $0.371$\,\scriptsize(+58.37\%) \\
\Cfive{} & $0.804$\,\scriptsize(+248.41\%) & $1.261$\,\scriptsize(+8.46\%) & $0.437$\,\scriptsize(+21.56\%) & $0.402$\,\scriptsize(+16.58\%) & $0.474$\,\scriptsize(-0.71\%) & $1.718$\,\scriptsize(+9.12\%) & $120.603$\,\scriptsize(+1.58\%) & $1.164$\,\scriptsize(+11.96\%) & $0.373$\,\scriptsize(+59.43\%) \\
\Csix{} & $0.270$\,\scriptsize(+16.97\%) & $1.249$\,\scriptsize(+7.43\%) & $0.457$\,\scriptsize(+27.11\%) & $0.349$\,\scriptsize(+1.22\%) & $0.502$\,\scriptsize(+5.20\%) & $1.643$\,\scriptsize(+4.33\%) & $120.523$\,\scriptsize(+1.51\%) & $1.078$\,\scriptsize(+3.68\%) & $0.270$\,\scriptsize(+15.55\%) \\
\bottomrule
\end{tabular}
\end{adjustbox}
\end{table*}

\begin{table*}[h]
\centering
\footnotesize
\setlength{\tabcolsep}{2.5pt}
\renewcommand{\arraystretch}{1.14}
\caption{Per-domain mean MSE on \TimeMMD{} for \TaTS{}.
Each cell is the mean over $4$ horizons $\times\, 3$ seeds (and $3$
backbones for trained methods).  $\Delta\%$ vs.\ \Cone{} in
parentheses.  Source: \texttt{summaries/per\_domain\_tats.csv}.}
\label{tab:perdom_tats}
\begin{adjustbox}{max width=\textwidth,center}
\begin{tabular}{lrrrrrrrrr}
\toprule
\textbf{Cond.} & \rotatebox{60}{\textbf{Agricul}} & \rotatebox{60}{\textbf{Climate}} & \rotatebox{60}{\textbf{Economy}} & \rotatebox{60}{\textbf{Energy}} & \rotatebox{60}{\textbf{Environ}} & \rotatebox{60}{\textbf{Health}} & \rotatebox{60}{\textbf{Securit}} & \rotatebox{60}{\textbf{SocialG}} & \rotatebox{60}{\textbf{Traffic}} \\
\midrule
\Cone{} & $0.162$ & $0.967$ & $0.127$ & $0.387$ & $0.290$ & $1.333$ & $115.282$ & $1.029$ & $0.189$ \\
\Cthree{} & $0.163$\,\scriptsize(+0.11\%) & $0.967$\,\scriptsize(-0.02\%) & $0.127$\,\scriptsize(-0.04\%) & $0.387$\,\scriptsize(+0.01\%) & $0.290$\,\scriptsize(-0.10\%) & $1.333$\,\scriptsize(-0.01\%) & $115.285$\,\scriptsize(+0.00\%) & $1.029$\,\scriptsize(+0.02\%) & $0.189$\,\scriptsize(-0.02\%) \\
\Cfour{} & $0.162$\,\scriptsize(+0.07\%) & $0.967$\,\scriptsize(-0.02\%) & $0.127$\,\scriptsize(-0.03\%) & $0.388$\,\scriptsize(+0.03\%) & $0.290$\,\scriptsize(-0.13\%) & $1.333$\,\scriptsize(-0.01\%) & $115.285$\,\scriptsize(+0.00\%) & $1.029$\,\scriptsize(-0.01\%) & $0.189$\,\scriptsize(-0.02\%) \\
\Ceight{} & $0.163$\,\scriptsize(+0.11\%) & $0.967$\,\scriptsize(-0.00\%) & $0.127$\,\scriptsize(-0.01\%) & $0.388$\,\scriptsize(+0.03\%) & $0.290$\,\scriptsize(-0.00\%) & $1.332$\,\scriptsize(-0.03\%) & $115.283$\,\scriptsize(+0.00\%) & $1.029$\,\scriptsize(-0.01\%) & $0.189$\,\scriptsize(+0.16\%) \\
\Cnine{} & $15.990$\,\scriptsize(+9749.88\%) & $3.317$\,\scriptsize(+242.95\%) & $0.426$\,\scriptsize(+235.06\%) & $1.435$\,\scriptsize(+270.30\%) & $0.407$\,\scriptsize(+40.37\%) & $2.009$\,\scriptsize(+50.73\%) & $125.877$\,\scriptsize(+9.19\%) & $3.290$\,\scriptsize(+219.74\%) & $4.212$\,\scriptsize(+2127.64\%) \\
\Ctwo{} & $15.990$\,\scriptsize(+9749.88\%) & $3.317$\,\scriptsize(+242.95\%) & $0.426$\,\scriptsize(+235.02\%) & $1.435$\,\scriptsize(+270.31\%) & $0.406$\,\scriptsize(+40.15\%) & $2.009$\,\scriptsize(+50.76\%) & $125.877$\,\scriptsize(+9.19\%) & $3.290$\,\scriptsize(+219.73\%) & $4.212$\,\scriptsize(+2127.65\%) \\
\Cfive{} & $15.990$\,\scriptsize(+9749.89\%) & $3.317$\,\scriptsize(+242.95\%) & $0.426$\,\scriptsize(+235.20\%) & $1.435$\,\scriptsize(+270.29\%) & $0.407$\,\scriptsize(+40.18\%) & $2.009$\,\scriptsize(+50.75\%) & $125.873$\,\scriptsize(+9.19\%) & $3.290$\,\scriptsize(+219.74\%) & $4.212$\,\scriptsize(+2127.64\%) \\
\Csix{} & $0.281$\,\scriptsize(+73.02\%) & $1.250$\,\scriptsize(+29.22\%) & $0.517$\,\scriptsize(+307.00\%) & $0.345$\,\scriptsize(-10.90\%) & $0.428$\,\scriptsize(+47.44\%) & $1.691$\,\scriptsize(+26.85\%) & $120.161$\,\scriptsize(+4.23\%) & $1.077$\,\scriptsize(+4.63\%) & $0.289$\,\scriptsize(+53.02\%) \\
\bottomrule
\end{tabular}
\end{adjustbox}
\end{table*}


\subsection{Per-horizon}
\label{sec:perhorizon_full}

The horizons available depend on the source domain's sampling
frequency: \textsc{Environment} (daily) uses
$\{48, 96, 192, 336\}$, \textsc{Health}/\textsc{Energy} (weekly) use
$\{12, 24, 36, 48\}$, and the remaining six monthly domains use
$\{6, 8, 10, 12\}$.  The table below merges across domains, so
columns with a particular horizon $H$ aggregate over only the
domains that include that horizon.

% AUTO-GENERATED by paper/figures/make_appendix_tables.py.
% Do not edit by hand; rerun the script after changing summaries.

\begin{table*}[h]
\centering
\footnotesize
\setlength{\tabcolsep}{2.5pt}
\renewcommand{\arraystretch}{1.14}
\caption{Per-horizon mean MSE for \Aurora{} on \TimeMMD{}.
Aggregated over domains, seeds, and (for trained methods)
backbones; $\Delta\%$ vs.\ \Cone{} at the same horizon in
parentheses.}
\label{tab:perhz_aurora}
\begin{adjustbox}{max width=\textwidth,center}
\begin{tabular}{lrrrrrrrrrr}
\toprule
\textbf{Cond.} & $H=6$ & $H=8$ & $H=10$ & $H=12$ & $H=24$ & $H=36$ & $H=48$ & $H=96$ & $H=192$ & $H=336$ \\
\midrule
\Cone{} & $11.556$ & $12.144$ & $12.866$ & $10.174$ & $0.897$ & $0.991$ & $0.841$ & $0.284$ & $0.270$ & $0.269$ \\
\Cthree{} & $11.557$\,\scriptsize(+0.01\%) & $12.145$\,\scriptsize(+0.01\%) & $12.866$\,\scriptsize(+0.00\%) & $10.173$\,\scriptsize(-0.00\%) & $0.897$\,\scriptsize(+0.00\%) & $0.991$\,\scriptsize(-0.01\%) & $0.841$\,\scriptsize(+0.00\%) & $0.284$\,\scriptsize(+0.00\%) & $0.270$\,\scriptsize(+0.00\%) & $0.269$\,\scriptsize(+0.00\%) \\
\Cfour{} & $11.563$\,\scriptsize(+0.06\%) & $12.147$\,\scriptsize(+0.02\%) & $12.860$\,\scriptsize(-0.04\%) & $10.169$\,\scriptsize(-0.05\%) & $0.897$\,\scriptsize(-0.01\%) & $0.990$\,\scriptsize(-0.06\%) & $0.841$\,\scriptsize(-0.03\%) & $0.284$\,\scriptsize(+0.00\%) & $0.270$\,\scriptsize(-0.00\%) & $0.269$\,\scriptsize(-0.00\%) \\
\Ceight{} & $11.573$\,\scriptsize(+0.14\%) & $12.152$\,\scriptsize(+0.07\%) & $12.862$\,\scriptsize(-0.03\%) & $10.165$\,\scriptsize(-0.09\%) & $0.899$\,\scriptsize(+0.22\%) & $0.994$\,\scriptsize(+0.33\%) & $0.843$\,\scriptsize(+0.28\%) & $0.284$\,\scriptsize(+0.00\%) & $0.270$\,\scriptsize(+0.00\%) & $0.269$\,\scriptsize(+0.01\%) \\
\Cnine{} & $11.556$\,\scriptsize(+0.00\%) & $12.144$\,\scriptsize(+0.00\%) & $12.866$\,\scriptsize(+0.00\%) & $10.174$\,\scriptsize(+0.00\%) & $0.897$\,\scriptsize(+0.00\%) & $0.991$\,\scriptsize(+0.00\%) & $0.841$\,\scriptsize(+0.00\%) & $0.284$\,\scriptsize(+0.00\%) & $0.270$\,\scriptsize(+0.00\%) & $0.269$\,\scriptsize(+0.00\%) \\
\Ctwo{} & $11.554$\,\scriptsize(-0.01\%) & $12.145$\,\scriptsize(+0.01\%) & $12.871$\,\scriptsize(+0.05\%) & $10.179$\,\scriptsize(+0.05\%) & $0.897$\,\scriptsize(+0.03\%) & $0.990$\,\scriptsize(-0.08\%) & $0.840$\,\scriptsize(-0.04\%) & $0.284$\,\scriptsize(-0.04\%) & $0.270$\,\scriptsize(-0.05\%) & $0.269$\,\scriptsize(-0.06\%) \\
\Cfive{} & $11.576$\,\scriptsize(+0.17\%) & $12.154$\,\scriptsize(+0.08\%) & $12.861$\,\scriptsize(-0.03\%) & $10.164$\,\scriptsize(-0.10\%) & $0.899$\,\scriptsize(+0.27\%) & $0.996$\,\scriptsize(+0.46\%) & $0.844$\,\scriptsize(+0.36\%) & $0.285$\,\scriptsize(+0.06\%) & $0.270$\,\scriptsize(+0.06\%) & $0.269$\,\scriptsize(+0.07\%) \\
\Csix{} & $11.556$\,\scriptsize(+0.00\%) & $12.144$\,\scriptsize(+0.00\%) & $12.866$\,\scriptsize(+0.00\%) & $10.174$\,\scriptsize(+0.00\%) & $0.897$\,\scriptsize(+0.00\%) & $0.991$\,\scriptsize(+0.00\%) & $0.841$\,\scriptsize(+0.00\%) & $0.284$\,\scriptsize(+0.00\%) & $0.270$\,\scriptsize(+0.00\%) & $0.269$\,\scriptsize(+0.00\%) \\
\bottomrule
\end{tabular}
\end{adjustbox}
\end{table*}


\begin{table*}[h]
\centering
\footnotesize
\setlength{\tabcolsep}{2.5pt}
\renewcommand{\arraystretch}{1.14}
\caption{Per-horizon mean MSE for \MMTSF{} on \TimeMMD{}.
Aggregated over domains, seeds, and (for trained methods)
backbones; $\Delta\%$ vs.\ \Cone{} at the same horizon in
parentheses.}
\label{tab:perhz_mmtsflib}
\begin{adjustbox}{max width=\textwidth,center}
\begin{tabular}{lrrrrrrrrrr}
\toprule
\textbf{Cond.} & $H=6$ & $H=8$ & $H=10$ & $H=12$ & $H=24$ & $H=36$ & $H=48$ & $H=96$ & $H=192$ & $H=336$ \\
\midrule
\Cone{} & $19.746$ & $20.173$ & $20.547$ & $15.708$ & $0.925$ & $1.041$ & $0.920$ & $0.486$ & $0.502$ & $0.471$ \\
\Cthree{} & $19.759$\,\scriptsize(+0.06\%) & $20.172$\,\scriptsize(-0.00\%) & $20.561$\,\scriptsize(+0.06\%) & $15.705$\,\scriptsize(-0.02\%) & $0.932$\,\scriptsize(+0.79\%) & $1.041$\,\scriptsize(+0.04\%) & $0.920$\,\scriptsize(+0.02\%) & $0.482$\,\scriptsize(-0.73\%) & $0.501$\,\scriptsize(-0.27\%) & $0.474$\,\scriptsize(+0.60\%) \\
\Cfour{} & $19.766$\,\scriptsize(+0.10\%) & $20.140$\,\scriptsize(-0.16\%) & $20.564$\,\scriptsize(+0.08\%) & $15.705$\,\scriptsize(-0.02\%) & $0.928$\,\scriptsize(+0.34\%) & $1.038$\,\scriptsize(-0.26\%) & $0.914$\,\scriptsize(-0.61\%) & $0.484$\,\scriptsize(-0.26\%) & $0.500$\,\scriptsize(-0.45\%) & $0.474$\,\scriptsize(+0.64\%) \\
\Ceight{} & $19.675$\,\scriptsize(-0.36\%) & $20.256$\,\scriptsize(+0.41\%) & $20.558$\,\scriptsize(+0.05\%) & $15.727$\,\scriptsize(+0.12\%) & $0.932$\,\scriptsize(+0.82\%) & $1.036$\,\scriptsize(-0.43\%) & $0.917$\,\scriptsize(-0.24\%) & $0.489$\,\scriptsize(+0.64\%) & $0.502$\,\scriptsize(+0.04\%) & $0.476$\,\scriptsize(+0.91\%) \\
\Cnine{} & $20.136$\,\scriptsize(+1.97\%) & $20.690$\,\scriptsize(+2.56\%) & $21.047$\,\scriptsize(+2.43\%) & $16.093$\,\scriptsize(+2.45\%) & $1.032$\,\scriptsize(+11.58\%) & $1.146$\,\scriptsize(+10.11\%) & $0.991$\,\scriptsize(+7.73\%) & $0.479$\,\scriptsize(-1.33\%) & $0.490$\,\scriptsize(-2.42\%) & $0.469$\,\scriptsize(-0.56\%) \\
\Ctwo{} & $20.133$\,\scriptsize(+1.96\%) & $20.721$\,\scriptsize(+2.71\%) & $21.019$\,\scriptsize(+2.30\%) & $16.107$\,\scriptsize(+2.54\%) & $1.032$\,\scriptsize(+11.67\%) & $1.157$\,\scriptsize(+11.22\%) & $0.986$\,\scriptsize(+7.23\%) & $0.479$\,\scriptsize(-1.30\%) & $0.488$\,\scriptsize(-2.73\%) & $0.470$\,\scriptsize(-0.35\%) \\
\Cfive{} & $20.127$\,\scriptsize(+1.93\%) & $20.733$\,\scriptsize(+2.77\%) & $21.020$\,\scriptsize(+2.30\%) & $16.111$\,\scriptsize(+2.57\%) & $1.034$\,\scriptsize(+11.81\%) & $1.156$\,\scriptsize(+11.06\%) & $0.988$\,\scriptsize(+7.48\%) & $0.481$\,\scriptsize(-0.95\%) & $0.488$\,\scriptsize(-2.71\%) & $0.469$\,\scriptsize(-0.44\%) \\
\Csix{} & $20.101$\,\scriptsize(+1.79\%) & $20.550$\,\scriptsize(+1.87\%) & $20.841$\,\scriptsize(+1.43\%) & $15.991$\,\scriptsize(+1.80\%) & $0.958$\,\scriptsize(+3.60\%) & $1.087$\,\scriptsize(+4.50\%) & $0.954$\,\scriptsize(+3.75\%) & $0.514$\,\scriptsize(+5.92\%) & $0.533$\,\scriptsize(+6.10\%) & $0.496$\,\scriptsize(+5.33\%) \\
\bottomrule
\end{tabular}
\end{adjustbox}
\end{table*}


\begin{table*}[h]
\centering
\footnotesize
\setlength{\tabcolsep}{2.5pt}
\renewcommand{\arraystretch}{1.14}
\caption{Per-horizon mean MSE for \TaTS{} on \TimeMMD{}.
Aggregated over domains, seeds, and (for trained methods)
backbones; $\Delta\%$ vs.\ \Cone{} at the same horizon in
parentheses.}
\label{tab:perhz_tats}
\begin{adjustbox}{max width=\textwidth,center}
\begin{tabular}{lrrrrrrrrrr}
\toprule
\textbf{Cond.} & $H=6$ & $H=8$ & $H=10$ & $H=12$ & $H=24$ & $H=36$ & $H=48$ & $H=96$ & $H=192$ & $H=336$ \\
\midrule
\Cone{} & $19.256$ & $19.494$ & $19.883$ & $15.081$ & $0.843$ & $0.906$ & $0.747$ & $0.298$ & $0.302$ & $0.285$ \\
\Cthree{} & $19.257$\,\scriptsize(+0.00\%) & $19.495$\,\scriptsize(+0.01\%) & $19.883$\,\scriptsize(+0.00\%) & $15.081$\,\scriptsize(+0.00\%) & $0.843$\,\scriptsize(+0.00\%) & $0.905$\,\scriptsize(-0.05\%) & $0.747$\,\scriptsize(+0.01\%) & $0.296$\,\scriptsize(-0.55\%) & $0.302$\,\scriptsize(+0.10\%) & $0.285$\,\scriptsize(-0.00\%) \\
\Cfour{} & $19.257$\,\scriptsize(+0.00\%) & $19.495$\,\scriptsize(+0.00\%) & $19.883$\,\scriptsize(+0.00\%) & $15.082$\,\scriptsize(+0.00\%) & $0.843$\,\scriptsize(+0.00\%) & $0.905$\,\scriptsize(-0.08\%) & $0.746$\,\scriptsize(-0.02\%) & $0.296$\,\scriptsize(-0.54\%) & $0.302$\,\scriptsize(+0.18\%) & $0.285$\,\scriptsize(+0.02\%) \\
\Ceight{} & $19.256$\,\scriptsize(-0.00\%) & $19.495$\,\scriptsize(+0.01\%) & $19.882$\,\scriptsize(-0.00\%) & $15.081$\,\scriptsize(+0.00\%) & $0.843$\,\scriptsize(+0.01\%) & $0.905$\,\scriptsize(-0.06\%) & $0.747$\,\scriptsize(+0.01\%) & $0.297$\,\scriptsize(-0.24\%) & $0.302$\,\scriptsize(+0.04\%) & $0.285$\,\scriptsize(+0.01\%) \\
\Cnine{} & $25.118$\,\scriptsize(+30.44\%) & $25.516$\,\scriptsize(+30.89\%) & $25.670$\,\scriptsize(+29.11\%) & $19.717$\,\scriptsize(+30.74\%) & $1.678$\,\scriptsize(+99.13\%) & $1.776$\,\scriptsize(+96.13\%) & $1.393$\,\scriptsize(+86.63\%) & $0.430$\,\scriptsize(+44.30\%) & $0.384$\,\scriptsize(+27.04\%) & $0.389$\,\scriptsize(+36.24\%) \\
\Ctwo{} & $25.118$\,\scriptsize(+30.44\%) & $25.516$\,\scriptsize(+30.89\%) & $25.670$\,\scriptsize(+29.11\%) & $19.717$\,\scriptsize(+30.74\%) & $1.678$\,\scriptsize(+99.11\%) & $1.777$\,\scriptsize(+96.16\%) & $1.393$\,\scriptsize(+86.53\%) & $0.428$\,\scriptsize(+43.87\%) & $0.384$\,\scriptsize(+27.32\%) & $0.389$\,\scriptsize(+36.21\%) \\
\Cfive{} & $25.116$\,\scriptsize(+30.43\%) & $25.516$\,\scriptsize(+30.89\%) & $25.670$\,\scriptsize(+29.11\%) & $19.717$\,\scriptsize(+30.74\%) & $1.678$\,\scriptsize(+99.11\%) & $1.777$\,\scriptsize(+96.15\%) & $1.393$\,\scriptsize(+86.59\%) & $0.428$\,\scriptsize(+43.73\%) & $0.384$\,\scriptsize(+27.30\%) & $0.388$\,\scriptsize(+36.16\%) \\
\Csix{} & $20.004$\,\scriptsize(+3.88\%) & $20.419$\,\scriptsize(+4.74\%) & $20.849$\,\scriptsize(+4.86\%) & $16.035$\,\scriptsize(+6.33\%) & $0.983$\,\scriptsize(+16.63\%) & $1.101$\,\scriptsize(+21.53\%) & $0.903$\,\scriptsize(+20.95\%) & $0.435$\,\scriptsize(+46.21\%) & $0.507$\,\scriptsize(+67.93\%) & $0.428$\,\scriptsize(+50.14\%) \\
\bottomrule
\end{tabular}
\end{adjustbox}
\end{table*}



\end{document}
